\documentclass[12pt]{article}
\usepackage[utf8]{inputenc}
\usepackage[spanish]{babel}
\usepackage{geometry}
\usepackage{enumitem}
\usepackage{titling}
\usepackage{fancyhdr}

\geometry{a4paper, margin=1in}

\title{Práctica de Física Básica: Cinemática y Leyes de Newton}
\author{Física Básica - Tema 2 y 3}
\date{\today}

\pagestyle{fancy}
\fancyhf{}
\rhead{Práctica de Física}
\lhead{Temas 2 y 3}
\cfoot{\thepage}

\begin{document}

\maketitle

\section*{Introducción}
Esta práctica abarca los conceptos fundamentales de la \textbf{Cinemática} (Tema 2) y los fundamentos de la \textbf{Dinámica} (Tema 3: Leyes de Newton), basada en el material del libro de Física Básica.

\hr
\section*{I. SECCIÓN TEÓRICA (10 Preguntas)}

\begin{enumerate}[label=\arabic*.]
    \item \textbf{La rama de la física que se encarga de la descripción del movimiento, sin considerar las causas que lo producen, es:}
    \begin{enumerate}[label=\alph*)]
        \item Dinámica
        \item Estática
        \item Cinemática
        \item Mecánica Cuántica
    \end{enumerate}

    \item \textbf{¿Cómo se denomina al ente constituido por un punto de referencia, un conjunto de ejes coordenados y un reloj?}
    \begin{enumerate}[label=\alph*)]
        \item Partícula
        \item Trayectoria
        \item Posición
        \item Marco de referencia
    \end{enumerate}

    \item \textbf{La línea imaginaria que describe un objeto durante su movimiento se conoce como:}
    \begin{enumerate}[label=\alph*)]
        \item Desplazamiento
        \item Trayectoria
        \item Distancia
        \item Vector posición
    \end{enumerate}

    \item \textbf{¿Cuál de las siguientes afirmaciones sobre el desplazamiento y la distancia es correcta?}
    \begin{enumerate}[label=\alph*)]
        \item El desplazamiento siempre es mayor que la distancia.
        \item Son iguales solo si el movimiento es en línea recta y en un solo sentido.
        \item La distancia es una cantidad vectorial.
        \item El desplazamiento es la longitud de la trayectoria.
    \end{enumerate}

    \item \textbf{En el Movimiento Rectilíneo Uniforme (MRU), ¿cuál de las siguientes características se mantiene constante?}
    \begin{enumerate}[label=\alph*)]
        \item La aceleración
        \item La posición
        \item La velocidad (magnitud, dirección y sentido)
        \item El tiempo
    \end{enumerate}

    \item \textbf{La aceleración se define fundamentalmente como el cambio en:}
    \begin{enumerate}[label=\alph*)]
        \item La posición por unidad de tiempo
        \item La velocidad por unidad de tiempo
        \item La trayectoria del objeto
        \item La masa del cuerpo
    \end{enumerate}

    \item \textbf{En el estudio de la Caída Libre, despreciando la resistencia del aire, la aceleración gravitacional (g) cerca de la superficie terrestre tiene una magnitud aproximada de:}
    \begin{enumerate}[label=\alph*)]
        \item 8.9 $m/s^2$ hacia arriba
        \item 9.8 $m/s^2$ hacia abajo
        \item 0 $m/s^2$
        \item 1.63 $m/s^2$ hacia abajo
    \end{enumerate}

    \item \textbf{La propiedad de los cuerpos de resistirse a los cambios en su estado de reposo o movimiento se denomina:}
    \begin{enumerate}[label=\alph*)]
        \item Fuerza
        \item Aceleración
        \item Inercia
        \item Velocidad
    \end{enumerate}

    \item \textbf{Según la Segunda Ley de Newton, la aceleración de un cuerpo es:}
    \begin{enumerate}[label=\alph*)]
        \item Inversamente proporcional a la fuerza neta
        \item Directamente proporcional a la masa
        \item Directamente proporcional a la fuerza neta e inversamente proporcional a la masa
        \item Independiente de la masa y la fuerza
    \end{enumerate}

    \item \textbf{La Tercera Ley de Newton establece que las fuerzas de acción y reacción:}
    \begin{enumerate}[label=\alph*)]
        \item Actúan sobre el mismo cuerpo y se anulan
        \item Tienen diferentes magnitudes
        \item Actúan sobre cuerpos diferentes y tienen sentidos contrarios
        \item Solo existen en el vacío
    \end{enumerate}
\end{enumerate}

\newpage
\section*{II. SECCIÓN PRÁCTICA (10 Ejercicios)}

\begin{enumerate}[label=\arabic*.]
    \item \textbf{Un vehículo se desplaza a una velocidad constante de 72 km/h. ¿Cuál es su velocidad expresada en metros por segundo (m/s)?}
    \begin{enumerate}[label=\alph*)]
        \item 10 m/s
        \item 20 m/s
        \item 25 m/s
        \item 720 m/s
    \end{enumerate}

    \item \textbf{Usain Bolt recorre 100 metros planos en un tiempo de 10.0 segundos. ¿Cuál fue su velocidad media durante la carrera?}
    \begin{enumerate}[label=\alph*)]
        \item 100 m/s
        \item 1.0 m/s
        \item 10 m/s
        \item 1000 m/s
    \end{enumerate}

    \item \textbf{Un tren viaja en línea recta con velocidad constante de 15 m/s durante 2 minutos. ¿Cuánto se desplazó el tren?}
    \begin{enumerate}[label=\alph*)]
        \item 30 m
        \item 150 m
        \item 1800 m
        \item 3000 m
    \end{enumerate}

    \item \textbf{Un auto que parte del reposo alcanza una velocidad de 20 m/s en un intervalo de 5.0 segundos. ¿Cuál fue su aceleración media?}
    \begin{enumerate}[label=\alph*)]
        \item 4.0 $m/s^2$
        \item 5.0 $m/s^2$
        \item 100 $m/s^2$
        \item 10 $m/s^2$
    \end{enumerate}

    \item \textbf{Determine el desplazamiento de un objeto que se mueve con MRUV si su velocidad inicial es de 5.0 m/s y acelera a 2.0 $m/s^2$ durante 4.0 segundos.}
    \begin{enumerate}[label=\alph*)]
        \item 20 m
        \item 16 m
        \item 36 m
        \item 40 m
    \end{enumerate}

    \item \textbf{Se deja caer una piedra desde lo alto de un edificio y tarda 3.0 segundos en tocar el suelo. ¿Cuál es la altura aproximada del edificio (g = 9.8 $m/s^2$)?}
    \begin{enumerate}[label=\alph*)]
        \item 29.4 m
        \item 44.1 m
        \item 88.2 m
        \item 14.7 m
    \end{enumerate}

    \item \textbf{Un proyectil es lanzado hacia arriba con una velocidad inicial de 20 m/s. ¿Qué velocidad tendrá al alcanzar su punto de máxima altura?}
    \begin{enumerate}[label=\alph*)]
        \item 20 m/s
        \item 9.8 m/s
        \item 0 m/s
        \item -9.8 m/s
    \end{enumerate}

    \item \textbf{¿Qué fuerza neta se requiere para impartirle una aceleración de 3.0 $m/s^2$ a un bloque de 5.0 kg de masa?}
    \begin{enumerate}[label=\alph*)]
        \item 1.66 N
        \item 8.0 N
        \item 15.0 N
        \item 2.0 N
    \end{enumerate}

    \item \textbf{¿Cuál es el peso de un objeto cuya masa es de 10.0 kg en la superficie de la Tierra (g = 9.8 $m/s^2$)?}
    \begin{enumerate}[label=\alph*)]
        \item 10.0 N
        \item 1.02 N
        \item 98.0 N
        \item 9.8 N
    \end{enumerate}

    \item \textbf{Un hombre de 70.0 kg va sentado en un auto que acelera a 1.0 $m/s^2$. La magnitud de la fuerza que el asiento ejerce sobre el hombre es:}
    \begin{enumerate}[label=\alph*)]
        \item 70.0 N
        \item 0 N
        \item 686.0 N
        \item 1.0 N
    \end{enumerate}
\end{enumerate}

\newpage
\section*{CLAVE DE RESPUESTAS}

\begin{table}[h]
\centering
\begin{tabular}{|c|c|c|c|}
\hline
\textbf{Pregunta (Teoría)} & \textbf{Respuesta} & \textbf{Pregunta (Práctica)} & \textbf{Respuesta} \\ \hline
1 & c & 1 & b \\ \hline
2 & d & 2 & c \\ \hline
3 & b & 3 & c \\ \hline
4 & b & 4 & a \\ \hline
5 & c & 5 & c \\ \hline
6 & b & 6 & b \\ \hline
7 & b & 7 & c \\ \hline
8 & c & 8 & c \\ \hline
9 & c & 9 & c \\ \hline
10 & c & 10 & a \\ \hline
\end{tabular}
\end{table}

\end{document}
